\documentclass[a4paper]{article}
\usepackage[margin=.8in]{geometry}
\usepackage{amsmath}
\usepackage{enumitem}
\usepackage[cache=true]{minted}
\usepackage{amsfonts}
\usepackage{amssymb}
\usepackage{mathtools}
\input{../../macros}

\usepackage{fancyhdr}
\pagestyle{fancy}
\fancyhf{}
\lhead{Name: }
\chead{Graded quiz 3 -- Interpolation}
\rhead{\today}

\begin{document}

\pagestyle{empty}
\thispagestyle{fancy}

True or false? (unless otherwise specified.)

\begin{center}
    \textbf{Answer grid}\par\smallskip

    \smallskip
    \normalsize
    \renewcommand{\arraystretch}{1.35}
    \begin{tabular}{|c|*{16}{c|}}
        \hline
        & 1 & 2 & 3 & 4 & 5 & 6 & 7 & 8 & 9 & 10 & 11 & 12 & 13 & 14 & 15 & B \\
        \hline
        True  & $\square$ & $\square$ & $\square$ & $\square$ & $\square$ & $\square$ & $\square$ & $\square$ & $\square$ & $\square$ & $\square$ & $\square$ & $\square$ & $\square$ & $\square$ & $\square$ \\
        False & $\square$ & $\square$ & $\square$ & $\square$ & $\square$ & $\square$ & $\square$ & $\square$ & $\square$ & $\square$ & $\square$ & $\square$ & $\square$ & $\square$ & $\square$ & $\square$ \\
        \hline
    \end{tabular}
\end{center}

\begin{enumerate}

    \item
        Let $x_0<x_1<x_2$. If $p,q\in\mathcal P_2$ satisfy
        \[
            p(x_i)=y_i,\qquad q(x_i)=z_i,\qquad i=0,1,2,
        \]
        then $p+q$ interpolates the data $(x_i,y_i+z_i)$ for $i=0,1,2$.

    \item
        For all $x_0,x_1,y_0,y_1\in\mathbb R$ with $x_0\ne x_1$, there exists a unique
        $p\in\mathcal P_2$ such that
        \[
            p(x_0)=y_0,\qquad p(x_1)=y_1.
        \]

    \item
        If $x_0<\dotsb<x_n$ and $\ell_i\in\mathcal P_n$ is the $i$th Lagrange basis polynomial,
        then by definition
        \[
            \ell_i(x_j)=
            \begin{cases}
                1, & i=j,\\
                0, & i\ne j,
            \end{cases}
            \qquad i,j=0,\dotsc,n.
        \]


    \item
        If $p,q\in\mathcal P_n$ and $p(x_i)=q(x_i)$ for $i=0,\dotsc,n$, where
        $x_0<\dotsb<x_n$, then $p=q$.

    \item
        If $f\in\mathcal P_d$ and $p_n\in\mathcal P_n$ is its interpolant at the $n+1$
        distinct nodes $x_0,\dotsc,x_n$, with $n\geq d$, then $p_n=f$.

    \item
        \textbf{Interpolation error.}
        Let $f\in C^{n+3}[a,b]$, let $a\leq x_0<\dotsb<x_n\leq b$, and let
        $p_n\in\mathcal P_n$ interpolate $f$ at $x_0,\dotsc,x_n$. Then, for every
        $x\in[a,b]$, there exists $\xi_x\in(a,b)$ such that
        \[
            f(x)-p_n(x)=
            \frac{f^{(n+3)}(\xi_x)}{(n+3)!}\prod_{i=0}^{n}(x-x_i).
        \]

    \item
        \textbf{Interpolation error.}
        Let $u(x)=x^4+2x^3+3x^2+4$ and let $p_3\in\mathcal P_3$ be its interpolant at
        the points $1,2,3,4$. Then
        \[
            u(x)-p_3(x)=(x-1)(x-2)(x-3)(x-4),\qquad \forall x\in\mathbb R.
        \]

    \item
        \textbf{Convergence.}
        For every function $f\in C[-1,1]$, if $p_n$ is the polynomial interpolating
        $f$ at the $n+1$ equally spaced nodes in $[-1,1]$, including $-1$ and $1$, then
        \[
            \lim_{n\to\infty} \left( \max_{x\in[-1,1]}\bigl\lvert p_n(x)-f(x)\bigr\rvert \right)=0.
        \]

    \item
        \textbf{Convergence.}
        Let $p_n \in \mathcal P_n$ be the polynomial interpolating
        $f(x)=\sin(2x)$ at the nodes $x_i=i/n$, for $i=0,\dotsc,n$. Then
        \[
            \lim_{n\to\infty} \left( \max_{x\in[0,1]} \bigl\lvert p_n(x)-f(x)\bigr\rvert \right) = + \infty.
        \]

    \item
        \textbf{Convergence estimate.}
        If $f\in C[a,b]$ and $f_n$ is the \textbf{piecewise quadratic} interpolant of $f$ based on a partition of $[a,b]$ into $n$ 
        equally large subintervals, then
        \[
            \lim_{n\to\infty}\left( \max_{a\leq x\leq b}\bigl\lvert f(x)-f_n(x)\bigr\rvert \right)=0.
        \]

    \item
        \textbf{Chebyshev polynomials.}
        Fix $n \in \mathbb N$ and let $T_n \colon[-1,1]\to\mathbb R$ be defined by $T_n(x) = \cos(n \arccos x)$.
        Then $T_n$ is a polynomial of degree $n$.

    \item
        \textbf{Chebyshev polynomials.}
        Assume that~$T_2$ is defined as in the previous item.
        Then $T_2(x) = 2x^2 - 3$.
        It may be useful to recall that~$\cos(2x) = 2\cos(x)^2 - 1$.

    \item
        \textbf{Chebyshev polynomials.}
        Let $x_0 < \dotsb < x_n$ denote the roots of the Chebyshev polynomial $T_{n+1}$ on $[-1,1]$.
        Using these points as interpolation nodes instead of equally spaced nodes can help reduce the maximum interpolation error.



    \item
        \textbf{Julia.}
        Let \julia{x = [0, 1, 2]} and \julia{y = [1, 2, 5]}. The following Julia code uses
        the Gregory--Newton approach to compute the interpolating polynomial through the three data
        points $(0,1)$, $(1,2)$, and $(2,5)$:
        \begin{minted}{julia}
function p(z)
    return y[1] + (y[2] - y[1]) * z
           + 1/2 * (y[3] - 2y[2] + y[1]) * z * (z + 1)
end
        \end{minted}

    \item
        \textbf{Julia.}
        Let \julia{x = [0, 1, 2]} and \julia{y = [1, 2, 5]}. The following
        Julia code correctly evaluates the Lagrange interpolating polynomial through these three
        data points at a point \julia{z}:
        \begin{minted}{julia}
function p(z)
    return (y[1] * (z - x[2]) * (z - x[3]) / ((x[1] - x[2]) * (x[1] - x[3]))
            + y[2] * (z - x[1]) * (z - x[3]) / ((x[2] - x[1]) * (x[2] - x[3]))
            + y[3] * (z - x[1]) * (z - x[2]) / ((x[3] - x[1]) * (x[3] - x[2])))
end
        \end{minted}

    \item \textbf{Bonus.}
        Let $f\in C^{\infty}[a,b]$, and let $f_n$ be the \textbf{piecewise linear} interpolant of $f$ on a
        uniform mesh with mesh size $h_n= \frac{b-a}{n}$. 
        Then
        \[
            \lim_{n \to \infty} n \times \left( \max_{x\in[a,b]}\bigl\lvert f(x)-f_n(x)\bigr\rvert \right)=0.
        \]
        \emph{Justify briefly.}

\end{enumerate}

\ifdefined\showsolutions
\newpage
\lhead{Instructor copy}
\section*{Answer key}
\noindent Questions 1--15: 1 point for a correct answer, 0 otherwise.

\begin{enumerate}[itemsep=8pt]
    \item \textbf{True.} For $i=0,1,2$, $(p+q)(x_i)=p(x_i)+q(x_i)=y_i+z_i$.
    \item \textbf{False.} Two points do not determine a unique polynomial in $\mathcal P_2$.
    \item \textbf{True.} This is the defining cardinal property of the Lagrange basis.
    \item \textbf{True.} The difference $p-q$ has degree at most $n$ and $n+1$ distinct roots.
    \item \textbf{True.} The difference $p_n-f$ belongs to $\mathcal P_n$ and has $n+1$ distinct roots.
    \item \textbf{False.} The interpolation remainder uses $f^{(n+1)}(\xi_x)/(n+1)!$, not
    $f^{(n+3)}(\xi_x)/(n+3)!$.
    \item \textbf{True.} Since $u^{(4)}=24$, the interpolation remainder is exactly the product of the four node factors.
    \item \textbf{False.} Equidistant interpolation need not converge uniformly for every continuous function.
    \item \textbf{False.} The function $\sin(2x)$ is analytic, and these interpolants converge uniformly on $[0,1]$, so the limit is $0$, not $+\infty$.
    \item \textbf{True.} Piecewise polynomial interpolation on a refining uniform mesh converges uniformly for continuous $f$.
    \item \textbf{True.} For $n\in\mathbb N$, $T_n$ is a polynomial of degree $n$.
    \item \textbf{False.} The identity is $T_2(x)=2x^2-1$, not $2x^2-3$.
    \item \textbf{True.} Chebyshev nodes generally reduce the maximum interpolation error compared with equally spaced nodes.
    \item \textbf{False.} The Gregory--Newton factor should be $z(z-1)$, not $z(z+1)$.
    \item \textbf{True.} The three displayed terms are the Lagrange basis contributions for the three nodes.
\end{enumerate}

\noindent\textbf{Bonus.}
\begin{enumerate}[itemsep=8pt]
    \item On each mesh interval $[x_i,x_{i+1}]$, the interpolation remainder gives
    \[
        f(x)-f_n(x)=\frac{f''(\xi_x)}{2}(x-x_i)(x-x_{i+1})
    \]
    for some $\xi_x\in(x_i,x_{i+1})$. Since
    \[
        \max_{x\in[x_i,x_{i+1}]}\lvert(x-x_i)(x-x_{i+1})\rvert
        =\frac{h_n^2}{4},
    \]
    the stated bound follows. Because $h_n\to0$, the right-hand side tends to zero.
\end{enumerate}
\fi

\end{document}
